% =============================================================================
% Section 2: Setup and notation
% Stubbed; populated as needed for §3 + downstream sections.
% =============================================================================

\section{Setup and notation}
\label{sec:setup}

\subsection{The divisor function and the framework constants}

\begin{notation}
We write $\tauf(n) = \#\{d \in \N : d \mid n\}$ for the number of
divisors of $n$, and $v_p(n)$ for the $p$-adic valuation of $n$. The
modulus
\[
M \;=\; 11 \cdot 13 \cdot 17 \cdot 19 \;=\; 46189
\]
is the squarefree composite of the four smallest primes that do
\emph{not} divide $\bigA / M = 2520 = 2^3 \cdot 3^2 \cdot 5 \cdot 7$.
We set $\Acap = 2520 \cdot M$ throughout.

The Stage-1 sieve (cf.\ \cite{ErdosGraham1980,GuyB8}) reduces the
$\tauf$-barrier search to $|\Ropen| = 96$ residue classes modulo $M$,
of which 55 have additional Lean-verified per-residue closure
certificates produced to date, with 41 remaining as the operational
frontier (cf.\ \cref{rem:5541-not-structural} on the artifact nature
of the $55/41$ split). We label this 41-element subset
$\Ropen \subset \Z/M\Z$ throughout.
\end{notation}

\subsection{The 143-collapse normalization}

\begin{notation}
By \cref{thm:143-collapse} every $r \in \Ropen$ satisfies
$143 \mid r$. We work in the equivalent normalized variables
\[
n \;=\; 143 \cdot Y , \qquad Y \in \N , \qquad Y \bmod 323 \in \Uset ,
\]
where $\Uset = \{r/143 \bmod 323 : r \in \Ropen\}$ is the explicit
41-element subset of $\Z/323\Z$ enumerated in \cref{thm:143-collapse}.
The shift form is
\[
2520 \cdot n - k \;=\; 360360 \cdot Y - k , \qquad
360360 \;=\; \lcm(1, 2, \ldots, 15) ,
\]
and we abbreviate $F_k(s) = 2520(Ms + r) - k = \Acap s + 2520 r - k$ for
$s \in \Z$ and $k \ge 1$.
\end{notation}

\subsection{Hensel-balanced kernel structure}

\begin{notation}
A \emph{Hensel-balanced kernel} for residue $r \in \Ropen$ is a CRT
class $s \equiv a \pmod{\Lcap}$ such that for every $1 \le k \le K$:
\[
F_k(s) \;=\; D_k(a) \cdot Q_k(t) , \qquad
Q_k(t) = \alpha_k + \beta_k \cdot t \in \Z[t] , \qquad
s = a + \Lcap \cdot t,
\]
where $D_k(a)$ is the forced divisor at $a$ and $\beta_k > 0$. The
kernel is \emph{admissible} if for every $k \le K_0$ (the
\emph{protected window}, typically $K_0 = 95$ or $K_0 = 1021$)
\[
2 \cdot \tauf(D_k(a)) \;\le\; k + 2 \quad \text{(Conjecture B)} .
\]

The kernel modulus $\Lcap$ is constructed as $\prod_p p^{E_p}$ for an
explicit cap structure; the cap formula is verified by direct enumeration in our internal Hensel-admissibility audit.
\end{notation}

\subsection{Pointwise-local-existential primitives}

The class \PLEclass\ of pointwise-local-existential forced-divisor
methods is defined in \cref{def:ple}. We additionally use the
following:

\begin{notation}
For a kernel $(a, \Lcap)$ and an auxiliary prime $\ell \le B$,
the \emph{safe set} at $\ell$ is
\[
S_\ell \;=\; \bigl\{ t \bmod \ell \,:\, Q_k(t) \not\equiv 0 \pmod \ell
\text{ for all } 1 \le k \le K_0 \bigr\} \,.
\]
A \emph{full safe-set CSP} chooses $a$ subject to admissibility and
$y_\ell \in S_\ell$ for each aux prime, and verifies that no shift
$k \in (K_0, K_{\rm scan}]$ produces a forced-divisor kill.
\end{notation}

\subsection{Divisor sieve and analytic NT references}

We use standard notation for sieves
(Iwaniec--Kowalski \cite{IwaniecKowalski2004}), Bombieri--Vinogradov,
the Bateman--Horn singular series \cite{BatemanHorn1962}, and
Schinzel's Hypothesis~H \cite{SchinzelSierpinski1958}. The
multidimensional Selberg weight of Maynard \cite{Maynard2013} and the
short-interval correlation framework of
Tao--Ter\"av\"ainen \cite{TaoTeravainen2025} and
Matom\"aki--Radziwi\l{}\l --Tao \cite{MRT2019} are reviewed in
\cref{sec:obstructions} for the obstruction analysis. The Hardy--
Tenenbaum normal order $\tauf(n) = (\log n)^{\log 2 + o(1)}$
\cite{HardyTenenbaum2008} is used in the analytic K-effective
calculations.

% =============================================================================
% End of section 2
% =============================================================================
